% !TEX root =esercizi.tex
% ---------------------------------------------------------------------
% Scheda 4: Lezioni 10--12. In L10 esponenziali e logaritmi come funzioni
% (completati) e funzioni trigonometriche; in L11 equazioni e disequazioni
% trigonometriche, inverse, campi di esistenza di riepilogo; in L12
% successioni, Fibonacci e rapporto aureo, successioni divergenti.
% Non svolti in L12 (e quindi senza esercizi qui): ricorrenza affine,
% monotonia e limitatezza (riprese in L13).
% ---------------------------------------------------------------------
\chapter{Lezioni 10--12 (14--19 ottobre 2026)}

\begin{vuoto}[Argomenti svolti]{}
\begin{itemize}
\item[\bul] \textbf{Lezione 10 --- Esponenziali e logaritmi; funzioni trigonometriche}. Funzioni esponenziali e logaritmiche: grafici, monotonia, crescita e decadimento, tempo di raddoppio e di dimezzamento, il numero $e$ e il logaritmo naturale, $N(t)=N_0e^{kt}$. Radianti; seno, coseno, tangente sulla circonferenza goniometrica; grafici.
\item[\bul] \textbf{Lezione 11 --- Trigonometria}. Valori notevoli e simmetrie; equazioni e disequazioni trigonometriche elementari, lette sulla circonferenza e sul grafico; funzioni trigonometriche inverse; modelli periodici $m+A\sin(\omega t+\varphi)$. Campi di esistenza di funzioni composte, come riepilogo di tutte le disequazioni.
\item[\bul] \textbf{Lezione 12 --- Successioni}. Successioni come funzioni su $\N$: elenco, formula esplicita, legge ricorsiva, grafico. La successione di Fibonacci, i rapporti fra termini consecutivi, il rapporto aureo e la sezione aurea. Successioni divergenti e confronto delle loro velocità.
\end{itemize}
\end{vuoto}

Il livello di difficoltà è indicato con \liv{1} (esercizio di base), \liv{2} (esercizio standard), \liv{3} (esercizio più impegnativo).

%-------------------------------------------------------------------
\section{Crescita e decadimento esponenziale}

\begin{esercizio}[\liv{1}]
Lo iodio-131, usato in medicina nucleare, ha un tempo di dimezzamento di circa $8$ giorni. Partendo da $100$ mg, quanto ne resta dopo $24$ giorni? E dopo $30$ giorni? Dopo quanto tempo ne resta $1$ mg?
\end{esercizio}
\risp{$N(t)=100\cdot2^{-t/8}$. Dopo $24$ giorni (tre dimezzamenti) $12{,}5$ mg; dopo $30$ giorni $100\cdot2^{-3{,}75}\approx7{,}4$ mg. $2^{-t/8}=0{,}01$ dà $t=8\log_2100\approx53$ giorni.}

\begin{esercizio}[\liv{2}]
Una popolazione cresce del $3\,\%$ all'anno. Calcolare il tempo di raddoppio esatto e quello dato dalla regola del $70$. Scrivere il modello nella forma $N_0e^{kt}$: quanto vale il tasso istantaneo $k$?
\end{esercizio}
\risp{$T=\dfrac{\log2}{\log1{,}03}\approx23{,}4$ anni; regola del $70$: $\frac{70}3\approx23{,}3$ anni. $k=\log1{,}03\approx0{,}0296$, cioè circa $2{,}96\,\%$, un po' meno del $3\,\%$.}

\begin{esercizio}[\liv{2}]
Il carbonio-14 ha tempo di dimezzamento di $5730$ anni. Un reperto contiene il $40\,\%$ del carbonio-14 che conteneva in origine. Quanti anni ha?
\end{esercizio}
\risp{$2^{-t/5730}=0{,}4$: $t=5730\log_2\frac1{0{,}4}=5730\,\dfrac{\log2{,}5}{\log2}\approx7{,}6\cdot10^3$ anni.}

\begin{esercizio}[\liv{2}]
Scrivere $N(t)=200\cdot0{,}85^{t}$ nella forma $N_0e^{kt}$. Si tratta di crescita o di decadimento? Qual è il tempo di dimezzamento?
\end{esercizio}
\risp{$k=\log0{,}85\approx-0{,}163$: decadimento, $N(t)=200\,e^{-0{,}163\,t}$. Il tempo di dimezzamento è $\dfrac{\log2}{-\log0{,}85}\approx4{,}27$.}

\begin{esercizio}[\liv{2}]
Il pH di una soluzione è $-\log_{10}[\mathrm H^+]$, con la concentrazione in mol/L. Calcolare il pH di una soluzione con $[\mathrm H^+]=3{,}2\cdot10^{-5}$ mol/L. Quante volte è più grande la concentrazione di ioni idrogeno in una soluzione a pH $4$ rispetto a una a pH $6$?
\end{esercizio}
\risp{$\mathrm{pH}=5-\log_{10}3{,}2\approx4{,}49$. A pH $4$ la concentrazione è $10^{-4}$, a pH $6$ è $10^{-6}$: è $100$ volte più grande.}

\begin{esercizio}[\liv{3}]
Due colture batteriche, con $t$ in ore: $A(t)=1000\cdot2^{t/3}$ e $B(t)=8000\cdot2^{t/5}$. Quale raddoppia più in fretta? Da quale istante in poi $A$ supera $B$?
\end{esercizio}
\risp{$A$ raddoppia in $3$ ore, $B$ in $5$. $A(t)>B(t)$ equivale a $2^{t/3-t/5}>8=2^3$, cioè $\frac{2t}{15}>3$: $t>22{,}5$ ore.}

%-------------------------------------------------------------------
\section{Radianti e valori delle funzioni trigonometriche}

\begin{esercizio}[\liv{1}]
Esprimere in radianti $30^\circ$, $135^\circ$, $270^\circ$, $-60^\circ$, e in gradi $\frac{5\pi}6$, $\frac{7\pi}4$, $2$ rad.
\end{esercizio}
\risp{$\frac\pi6$, $\frac{3\pi}4$, $\frac{3\pi}2$, $-\frac\pi3$; \ $150^\circ$, $315^\circ$, $\frac{360^\circ}{\pi}\approx114{,}6^\circ$.}

\begin{esercizio}[\liv{1}]
Calcolare: \ $\sin\frac{5\pi}6$; \ $\cos\frac{4\pi}3$; \ $\tan\frac{3\pi}4$; \ $\sin\bigl(-\frac\pi4\bigr)$; \ $\cos\frac{11\pi}6$; \ $\tan\frac{7\pi}6$; \ $\sin\frac{13\pi}3$.
\end{esercizio}
\risp{$\frac12$; \ $-\frac12$; \ $-1$; \ $-\frac{\sqrt2}2$; \ $\frac{\sqrt3}2$; \ $\frac{\sqrt3}3$; \ $\frac{13\pi}3=4\pi+\frac\pi3$, quindi $\frac{\sqrt3}2$.}

\begin{esercizio}[\liv{2}]
Sapendo che $\sin x=\frac35$ e $\frac\pi2<x<\pi$, calcolare $\cos x$, $\tan x$, $\sin(\pi-x)$, $\cos(x+\pi)$.
\end{esercizio}
\risp{Nel secondo quadrante il coseno è negativo: $\cos x=-\sqrt{1-\frac9{25}}=-\frac45$; $\tan x=-\frac34$; $\sin(\pi-x)=\sin x=\frac35$; $\cos(x+\pi)=-\cos x=\frac45$.}

\begin{esercizio}[\liv{2}]
Il raggio della Terra è circa $6370$ km. Quanto è lungo un arco di meridiano di $1^\circ$? Napoli si trova a una latitudine di circa $40{,}85^\circ$ nord: quanto dista dall'equatore, misurando lungo il meridiano?
\end{esercizio}
\risp{L'arco è lungo $r\theta$ con $\theta$ in radianti: $6370\cdot\frac\pi{180}\approx111$ km; $6370\cdot40{,}85\cdot\frac\pi{180}\approx4540$ km.}

%-------------------------------------------------------------------
\section{Grafici e modelli periodici}

\begin{esercizio}[\liv{1}]
Determinare periodo, ampiezza e immagine di: \ $3\sin(2x)$; \ $1+\cos\frac x2$; \ $1-2\sin(\pi x)$.
\end{esercizio}
\risp{Periodo $\pi$, ampiezza $3$, immagine $[-3,3]$; \ periodo $4\pi$, ampiezza $1$, immagine $[0,2]$; \ periodo $2$, ampiezza $2$, immagine $[-1,3]$.}

\begin{esercizio}[\liv{2}]
A partire dal grafico di $\sin x$ disegnare i grafici di $\sin\bigl(x-\frac\pi4\bigr)$, $2\cos x-1$, $|\sin x|$, $\sin|x|$. Quali sono periodici, e con quale periodo?
\end{esercizio}
\risp{$\sin(x-\frac\pi4)$: traslato a destra di $\frac\pi4$, periodo $2\pi$. $2\cos x-1$: dilatato verticalmente di $2$ e abbassato di $1$, periodo $2\pi$, immagine $[-3,1]$. $|\sin x|$: le parti negative ribaltate, periodo $\pi$. $\sin|x|$: per $x\ge0$ coincide con $\sin x$, ed è pari; non è periodica.}

\begin{esercizio}[\liv{2}]
In un porto l'altezza della marea, in metri, è descritta da $h(t)=1{,}2+0{,}8\sin\dfrac{2\pi t}{12{,}4}$, con $t$ in ore. Determinare altezza massima e minima e il periodo. In quale intervallo di tempo, fra $t=0$ e $t=12{,}4$, l'altezza supera $1{,}6$ m? Per quante ore?
\end{esercizio}
\risp{Massimo $2$ m, minimo $0{,}4$ m, periodo $12{,}4$ ore. $h>1{,}6$ equivale a $\sin\frac{2\pi t}{12{,}4}>\frac12$, cioè $\frac\pi6<\frac{2\pi t}{12{,}4}<\frac{5\pi}6$: $\frac{12{,}4}{12}<t<\frac{5\cdot12{,}4}{12}$, cioè $1{,}03<t<5{,}17$, per circa $4{,}1$ ore.}

\begin{esercizio}[\liv{3}]
In una località la temperatura media giornaliera ha un massimo di $26\,^\circ$C il giorno $200$ dell'anno e un minimo di $6\,^\circ$C. Scrivere un modello $T(t)=m+A\sin\bigl(\frac{2\pi}{365}(t-t_0)\bigr)$ e dire in quale giorno si ha il minimo.
\end{esercizio}
\risp{$m=\frac{26+6}2=16$, $A=\frac{26-6}2=10$. Il seno vale $1$ quando l'argomento è $\frac\pi2$, cioè $t-t_0=\frac{365}4$: $t_0=200-91{,}25=108{,}75$. $T(t)=16+10\sin\bigl(\frac{2\pi}{365}(t-108{,}75)\bigr)$. Il minimo è mezzo periodo dopo il massimo: $200+182{,}5-365=17{,}5$, cioè verso il 17--18 gennaio.}

%-------------------------------------------------------------------
\section{Equazioni trigonometriche}

\begin{esercizio}[\liv{1}]
Risolvere in $\R$: \ $\sin x=\frac{\sqrt3}2$; \ $\cos x=-\frac12$; \ $\tan x=1$; \ $\sin x=2$; \ $\cos x=0$.
\end{esercizio}
\risp{$x=\frac\pi3+2k\pi$ oppure $x=\frac{2\pi}3+2k\pi$; \ $x=\pm\frac{2\pi}3+2k\pi$; \ $x=\frac\pi4+k\pi$; \ nessuna soluzione; \ $x=\frac\pi2+k\pi$ ($k\in\Z$).}

\begin{esercizio}[\liv{2}]
Risolvere in $[0,2\pi]$: \ $\sin2x=\frac12$; \ $\cos\bigl(x-\frac\pi3\bigr)=1$; \ $2\sin^2x-\sin x-1=0$.
\end{esercizio}
\risp{$2x=\frac\pi6+2k\pi$ oppure $2x=\frac{5\pi}6+2k\pi$: $x=\frac\pi{12},\frac{5\pi}{12},\frac{13\pi}{12},\frac{17\pi}{12}$. \ $x-\frac\pi3=2k\pi$: $x=\frac\pi3$. \ Con $t=\sin x$: $2t^2-t-1=0$, $t=1$ o $t=-\frac12$: $x=\frac\pi2,\frac{7\pi}6,\frac{11\pi}6$.}

\begin{esercizio}[\liv{2}]
Risolvere in $\R$: \ $\sin x=\cos x$; \ $2\cos^2x+3\sin x-3=0$.
\end{esercizio}
\risp{Se $\cos x=0$ allora $\sin x=\pm1\ne0$, quindi si può dividere: $\tan x=1$, $x=\frac\pi4+k\pi$. \ Con $\cos^2x=1-\sin^2x$: $2\sin^2x-3\sin x+1=0$, $\sin x=1$ o $\sin x=\frac12$: $x=\frac\pi2+2k\pi$, $x=\frac\pi6+2k\pi$, $x=\frac{5\pi}6+2k\pi$.}

\begin{esercizio}[\liv{2}]
Risolvere $\sin x=0{,}3$ in $[0,2\pi]$, usando la funzione $\arcsin$, e dare i valori approssimati.
\end{esercizio}
\risp{$x=\arcsin0{,}3\approx0{,}305$ e $x=\pi-\arcsin0{,}3\approx2{,}837$.}

%-------------------------------------------------------------------
\section{Disequazioni trigonometriche}

\begin{esercizio}[\liv{1}]
Risolvere in $[0,2\pi]$: \ $\sin x\ge\frac12$; \ $\cos x<\frac{\sqrt2}2$; \ $\tan x>1$.
\end{esercizio}
\risp{$[\frac\pi6,\frac{5\pi}6]$; \ $]\frac\pi4,\frac{7\pi}4[$; \ $]\frac\pi4,\frac\pi2[\,\cup\,]\frac{5\pi}4,\frac{3\pi}2[$.}

\begin{esercizio}[\liv{2}]
Risolvere in $[0,2\pi]$: \ $2\cos x+1\le0$; \ $\sin^2x<\frac14$.
\end{esercizio}
\risp{$\cos x\le-\frac12$: $[\frac{2\pi}3,\frac{4\pi}3]$. \ $|\sin x|<\frac12$: $[0,\frac\pi6[\,\cup\,]\frac{5\pi}6,\frac{7\pi}6[\,\cup\,]\frac{11\pi}6,2\pi]$.}

\begin{esercizio}[\liv{2}]
Risolvere in $[0,2\pi]$ la disequazione $\sin2x>0$, ponendo $t=2x$ (attenzione all'intervallo in cui varia $t$).
\end{esercizio}
\risp{$t\in[0,4\pi]$ e $\sin t>0$ per $t\in\,]0,\pi[\,\cup\,]2\pi,3\pi[$: $x\in\,]0,\frac\pi2[\,\cup\,]\pi,\frac{3\pi}2[$.}

\begin{esercizio}[\liv{3}]
Risolvere in $[0,2\pi]$: \ $2\sin^2x-3\sin x+1>0$.
\end{esercizio}
\risp{Con $t=\sin x$: $(2t-1)(t-1)>0$, cioè $t<\frac12$ oppure $t>1$; la seconda è impossibile. Quindi $\sin x<\frac12$: $[0,\frac\pi6[\,\cup\,]\frac{5\pi}6,2\pi]$.}

%-------------------------------------------------------------------
\section{Funzioni trigonometriche inverse}

\begin{esercizio}[\liv{1}]
Calcolare: \ $\arcsin\frac12$; \ $\arccos\bigl(-\frac{\sqrt3}2\bigr)$; \ $\arctan(-1)$; \ $\arcsin(-1)$; \ $\arccos0$; \ $\arctan\sqrt3$.
\end{esercizio}
\risp{$\frac\pi6$; \ $\frac{5\pi}6$; \ $-\frac\pi4$; \ $-\frac\pi2$; \ $\frac\pi2$; \ $\frac\pi3$.}

\begin{esercizio}[\liv{2}]
Calcolare:
\[
\arcsin\Bigl(\sin\frac{3\pi}4\Bigr),\quad \arccos\Bigl(\cos\Bigl(-\frac\pi6\Bigr)\Bigr),\quad \arctan\Bigl(\tan\frac{5\pi}4\Bigr),\quad \sin\Bigl(\arccos\frac13\Bigr),\quad \tan\Bigl(\arcsin\frac35\Bigr).
\]
\end{esercizio}
\risp{$\frac\pi4$ (non $\frac{3\pi}4$, che non sta in $[-\frac\pi2,\frac\pi2]$); \ $\frac\pi6$; \ $\frac\pi4$; \ $\sqrt{1-\frac19}=\frac{2\sqrt2}3$ (positivo, perché $\arccos\frac13\in[0,\pi]$); \ con $\sin y=\frac35$ e $y\in[-\frac\pi2,\frac\pi2]$ si ha $\cos y=\frac45$, quindi $\frac34$.}

\begin{esercizio}[\liv{2}]
Determinare il dominio di: \ $\arcsin(2x-1)$; \ $\arccos\frac1x$; \ $\arcsin(x^2-1)$.
\end{esercizio}
\risp{$-1\le2x-1\le1$: $[0,1]$; \ $\bigl|\frac1x\bigr|\le1$: $|x|\ge1$; \ $0\le x^2\le2$: $[-\sqrt2,\sqrt2]$.}

\begin{esercizio}[\liv{2}]
Risolvere: \ $\arctan x>\frac\pi3$; \ $\arcsin x\le\frac\pi6$; \ $\arccos x<\frac\pi4$.
\end{esercizio}
\risp{$\arctan$ è crescente: $x>\sqrt3$. \ $\arcsin$ è crescente su $[-1,1]$: $[-1,\frac12]$. \ $\arccos$ è decrescente: $x>\frac{\sqrt2}2$, con $x\le1$: $]\frac{\sqrt2}2,1]$.}

%-------------------------------------------------------------------
\section{Campi di esistenza}

\begin{vuoto}[Promemoria]{}
Per determinare il dominio di una funzione composta si scrivono, e si mettono a sistema, le condizioni di ciascun ``pezzo'':
\begin{itemize}
\item[\bul] denominatori diversi da $0$;
\item[\bul] argomenti delle radici di indice pari maggiori o uguali a $0$;
\item[\bul] argomenti dei logaritmi strettamente positivi;
\item[\bul] argomenti di $\arcsin$ e $\arccos$ compresi fra $-1$ e $1$;
\item[\bul] argomenti della tangente diversi da $\frac\pi2+k\pi$.
\end{itemize}
\end{vuoto}

\begin{esercizio}[\liv{1}]
Determinare il dominio di: \ $\sqrt{x^2-5x+6}$; \ $\log(4-x^2)$; \ $\dfrac1{e^x-1}$; \ $\sqrt{\log x}$; \ $\dfrac{x+1}{\sqrt{9-x^2}}$.
\end{esercizio}
\risp{$x\le2$ oppure $x\ge3$; \ $]-2,2[$; \ $x\ne0$; \ $\log x\ge0$: $x\ge1$; \ $9-x^2>0$: $]-3,3[$.}

\begin{esercizio}[\liv{2}]
Determinare il dominio di: \ $\sqrt{\dfrac{x-1}{x+3}}$; \ $\log(|x|-2)$; \ $\sqrt{1-\log_2x}$; \ $\log_{1/3}(x^2-x)$; \ $\dfrac{\sqrt{2^x-8}}{x-5}$.
\end{esercizio}
\risp{$x<-3$ oppure $x\ge1$; \ $|x|>2$: $x<-2$ oppure $x>2$; \ $x>0$ e $\log_2x\le1$: $]0,2]$; \ $x^2-x>0$: $x<0$ oppure $x>1$; \ $2^x\ge8$ e $x\ne5$: $[3,5[\,\cup\,]5,+\infty[$.}

\begin{esercizio}[\liv{2}]
Determinare il dominio di:
\[
\sqrt{\log(x-2)},\qquad \log(\log x),\qquad \sqrt{\log_{1/2}(x+1)},\qquad \frac1{\sqrt{x+4}-2},\qquad \sqrt{4^x-3\cdot2^x+2}.
\]
\end{esercizio}
\risp{$x>2$ e $\log(x-2)\ge0$, cioè $x-2\ge1$: $[3,+\infty[$; \ $x>0$ e $\log x>0$: $x>1$; \ $x>-1$ e, con base minore di $1$, $x+1\le1$: $]-1,0]$; \ $x\ge-4$ e $\sqrt{x+4}\ne2$, cioè $x\ne0$: $[-4,0[\,\cup\,]0,+\infty[$; \ con $t=2^x>0$: $t^2-3t+2\ge0$, cioè $t\le1$ o $t\ge2$: $x\le0$ oppure $x\ge1$.}

\begin{esercizio}[\liv{2}]
Determinare il dominio di: \ $\sqrt{\sin x}$; \ $\log(\cos x)$; \ $\dfrac1{2\sin x-1}$; \ $\sqrt{1-2\cos x}$ (in $[0,2\pi]$); \ $\arcsin(\log x)$.
\end{esercizio}
\risp{$\sin x\ge0$: $[2k\pi,\pi+2k\pi]$; \ $\cos x>0$: $]-\frac\pi2+2k\pi,\frac\pi2+2k\pi[$; \ $\sin x\ne\frac12$: $x\ne\frac\pi6+2k\pi$ e $x\ne\frac{5\pi}6+2k\pi$; \ $\cos x\le\frac12$: $[\frac\pi3,\frac{5\pi}3]$; \ $x>0$ e $-1\le\log x\le1$: $[\frac1e,e]$ ($k\in\Z$).}

\begin{esercizio}[\liv{3}]
Determinare il dominio di: \ $\sqrt{(2-x)\log x}$; \ $\sqrt{\sqrt{x+3}-x-1}$; \ $\log\dfrac{x^2-4}{x-1}$; \ $\sqrt{\arctan x-\frac\pi4}$.
\end{esercizio}
\risp{Prima: $x>0$ e $(2-x)\log x\ge0$; i due fattori hanno lo stesso segno (o uno è nullo) solo per $1\le x\le2$: $[1,2]$. \ Seconda: $x\ge-3$ e $\sqrt{x+3}\ge x+1$; se $x<-1$ è vera; se $x\ge-1$ si eleva al quadrato: $x^2+x-2\le0$, $-2\le x\le1$. Soluzione $[-3,1]$. \ Terza: $\frac{(x-2)(x+2)}{x-1}>0$: $]-2,1[\,\cup\,]2,+\infty[$. \ Quarta: $\arctan x\ge\frac\pi4$: $x\ge1$.}

\begin{esercizio}[\liv{3}]
Determinare il dominio di $f(x)=\sqrt{\log_{1/2}(x^2-3x+3)}$.
\end{esercizio}
\risp{$x^2-3x+3>0$ per ogni $x$ (discriminante negativo). Con base minore di $1$, $\log_{1/2}(x^2-3x+3)\ge0$ equivale a $x^2-3x+3\le1$, cioè $x^2-3x+2\le0$: $[1,2]$.}

%-------------------------------------------------------------------
\section{Successioni}

\begin{esercizio}[\liv{1}]
Scrivere i primi cinque termini e disegnare il grafico di: \ $a_n=\dfrac{n-1}n$ ($n\ge1$); \ $b_n=(-1)^n\,n$ ($n\ge1$); \ $c_n=1+\dfrac{(-1)^n}n$ ($n\ge1$); \ $d_0=1$, $d_{n+1}=3d_n-1$.
\end{esercizio}
\risp{$0,\frac12,\frac23,\frac34,\frac45$; \ $-1,2,-3,4,-5$; \ $0,\frac32,\frac23,\frac54,\frac45$; \ $1,2,5,14,41$. I grafici sono insiemi di punti isolati, da non unire.}

\begin{esercizio}[\liv{1}]
Trovare una formula esplicita per una successione che inizia con: \ $5,8,11,14,\dots$; \ $3,-6,12,-24,\dots$; \ $\frac12,\frac23,\frac34,\frac45,\dots$; \ $1,8,27,64,\dots$
\end{esercizio}
\risp{$5+3n$ ($n\ge0$); \ $3(-2)^n$ ($n\ge0$); \ $\frac n{n+1}$ ($n\ge1$); \ $n^3$ ($n\ge1$). Non sono le uniche possibili: un numero finito di termini non determina la successione.}

\begin{esercizio}[\liv{2}]
Scrivere i numeri di Fibonacci fino a $F_{15}$ e calcolare i rapporti $F_{n+1}/F_n$ per $n=10,\dots,14$, confrontandoli con $\varphi=\frac{1+\sqrt5}2$. Verificare che $F_{n+1}F_{n-1}-F_n^2=(-1)^n$ per $n=2,\dots,6$.
\end{esercizio}
\risp{$1,1,2,3,5,8,13,21,34,55,89,144,233,377,610$. Rapporti: $1{,}618182$; $1{,}617978$; $1{,}618056$; $1{,}618026$; $1{,}618037$, alternativamente sopra e sotto $\varphi\approx1{,}618034$. Per esempio $F_3F_1-F_2^2=2-1=1$, $F_4F_2-F_3^2=3-4=-1$, $F_5F_3-F_4^2=10-9=1$.}

\begin{esercizio}[\liv{2}]
Un segmento di $10$ cm è diviso in sezione aurea: quanto sono lunghe le due parti? Un rettangolo aureo ha il lato minore di $5$ cm: quanto è lungo il lato maggiore? Verificare con la calcolatrice che $\varphi^2=\varphi+1$ e $\frac1\varphi=\varphi-1$.
\end{esercizio}
\risp{La parte maggiore è $\frac{10}\varphi\approx6{,}18$ cm, la minore $\frac{10}{\varphi^2}\approx3{,}82$ cm. Il lato maggiore è $5\varphi\approx8{,}09$ cm. $\varphi^2\approx2{,}618=\varphi+1$; $\frac1\varphi\approx0{,}618=\varphi-1$.}

\begin{esercizio}[\liv{2}]
Ordinare le successioni $n^3$, $2^n$, $\log n$, $n!$, $\sqrt n$, $3^n$, $n^n$ calcolandone i valori per $n=10$. L'ordine trovato è quello ``per $n$ grande''? Confrontare $n^3$ e $2^n$ per $n=9$ e $n=10$.
\end{esercizio}
\risp{Per $n=10$: $\log n\approx2{,}3<\sqrt n\approx3{,}2<n^3=1000<2^n=1024<3^n=59\,049<n!\approx3{,}6\cdot10^6<n^n=10^{10}$. È l'ordine per $n$ grande, ma per $n=9$ si ha ancora $n^3=729>2^n=512$: il sorpasso di $2^n$ su $n^3$ avviene proprio a $n=10$.}

\begin{esercizio}[\liv{3}]
Sia $a_1=a_2=1$ e $a_{n+1}=a_n+2a_{n-1}$. (a) Scrivere i primi dieci termini e i rapporti $r_n=a_{n+1}/a_n$. (b) Mostrare che $r_n=1+\dfrac2{r_{n-1}}$, e trovare il candidato per il numero a cui si avvicinano i rapporti, come per Fibonacci.
\end{esercizio}
\risp{(a) $1,1,3,5,11,21,43,85,171,341$; rapporti $1$; $3$; $1{,}667$; $2{,}2$; $1{,}909$; $2{,}048$; $1{,}977$; $2{,}012$; $1{,}994$. (b) Dividendo per $a_n$: $r_n=1+2\frac{a_{n-1}}{a_n}=1+\frac2{r_{n-1}}$. Il candidato soddisfa $x=1+\frac2x$, cioè $x^2-x-2=0$, e la radice positiva è $x=2$.}

\begin{esercizio}[\liv{3}]
Il perimetro del poligono regolare di $n$ lati inscritto in una circonferenza di diametro $1$ è $p_n=n\sin\frac\pi n$. Calcolare esattamente $p_4$, $p_6$ e, con la formula di sottrazione $\sin(\alpha-\beta)=\sin\alpha\cos\beta-\cos\alpha\sin\beta$ applicata a $\frac\pi{12}=\frac\pi3-\frac\pi4$, anche $p_{12}$. Confrontare con $\pi$.
\end{esercizio}
\risp{$p_4=4\cdot\frac{\sqrt2}2=2\sqrt2\approx2{,}83$; $p_6=6\cdot\frac12=3$; $\sin\frac\pi{12}=\frac{\sqrt3}2\cdot\frac{\sqrt2}2-\frac12\cdot\frac{\sqrt2}2=\frac{\sqrt6-\sqrt2}4$, quindi $p_{12}=3(\sqrt6-\sqrt2)\approx3{,}106$. Tutti minori di $\pi\approx3{,}1416$, e sempre più vicini.}

%-------------------------------------------------------------------
\stampasoluzioni
